\section*{Chapitre \ref{ch:ll:testing-and-deploying-low-latency-systems} --- Tests et déploiement des systèmes à faible latence}

\begin{solution}{exo:ll:testing-and-deploying-low-latency-systems:1}
Aucun ordre n'est jamais exécuté au-delà de sa quantité demandée~; la position ne dépasse jamais la limite comptée dans le pire cas (ou~: la
position est égale à la somme des exécutions). Que la vue de la passerelle concorde avec les registres de la plateforme demande un relevé indépendant
du côté de la plateforme, le \omterm{def:ll:order-gateway-and-order-management:dropcopy}{drop copy}.
\end{solution}

\begin{solution}{exo:ll:testing-and-deploying-low-latency-systems:2}
Ainsi, une lecture d'un octet au-delà de l'entrée est une lecture au-delà d'une allocation, ce que signale AddressSanitizer. Dans un tampon plus grand
(capacité d'une chaîne, tableau réutilisé), la même mauvaise lecture renverrait silencieusement des valeurs parasites.
\end{solution}

\begin{solution}{exo:ll:testing-and-deploying-low-latency-systems:3}
Tout ce qui varie avec le temps (le démarrage d'une autre tâche, la montée en température du processeur, un changement de fréquence) n'affecte alors qu'une construction et
non l'autre, et se mesure comme un écart entre les deux. L'entrelacement répartit ces dérives sur les deux.
\end{solution}

\begin{solution}{exo:ll:testing-and-deploying-low-latency-systems:4}
$2\,(1.645 + 1.282)^2\,(1.2533 \times 0.14 / r)^2$~: environ 53 passages par côté pour $r = 10~\%$, environ 840 pour $r = 2.5~\%$. Le
nombre croît comme l'inverse du carré de la hausse.
\end{solution}

\begin{solution}{exo:ll:testing-and-deploying-low-latency-systems:5}
$20 \times 21 = 420$ barrières~; à un seuil de 5~\% sans changement réel, jusqu'à 21 fausses alertes. La tolérance ne signale que les hausses
de plus de 2~\%, ce qui en supprime beaucoup~: le taux mesuré sur cet ordinateur portable était de 3 à 9~\%, soit 13 à 38 par mois~; une tolérance plus large ou
un seuil plus bas les réduirait, au prix de la puissance statistique. Chaque fausse alerte coûte une nouvelle mesure, si bien que le seuil et la tolérance se règlent en tenant compte de la patience de l'équipe.
\end{solution}

\begin{solution}{exo:ll:testing-and-deploying-low-latency-systems:6}
La somme de contrôle est la somme des octets modulo 256. Remplacer «~1~» par «~0~» diminue la somme d'une unité, remplacer «~2~» par «~3~» l'augmente
d'une unité~: la somme reste inchangée. Elle détecte la plupart des altérations d'un seul octet et les troncatures~; elle ne détecte pas les modifications compensatoires
ni les réordonnancements d'octets, ni toute altération délibérée. La protection contre la corruption sur le réseau relève de TCP~; les contrôles de l'analyseur
portent sur la structure.
\end{solution}

\begin{solution}{exo:ll:testing-and-deploying-low-latency-systems:7}
Une plateforme configurée avec un débit de messages (le limiteur du simulateur), une rafale d'ordres envoyés avec le même horodatage, et
des résultats attendus pour une suite d'acceptations suivies de rejets de motif \texttt{T}~; l'étape dépend du temps, si bien que le
lanceur de tests doit maîtriser les horodatages et la configuration du moteur, ce que les autres étapes n'exigeaient pas.
\end{solution}

\begin{solution}{exo:ll:testing-and-deploying-low-latency-systems:8}
Un passage de chaque, à une semaine d'intervalle~: la machine, sa charge, sa température et le reste du système diffèrent entre les deux, et
sur cet ordinateur portable, le 99\up{e}~percentile d'un seul passage varie de bien plus de 5~\%. La barrière échouerait au hasard et laisserait passer les vraies
régressions. Lancer les deux constructions maintenant, entrelacées, de nombreuses fois, et les comparer par un test.
\end{solution}

\begin{solution}{pb:ll:testing-and-deploying-low-latency-systems:1}
\begin{enumerate}
\item Environ \qty{208}{\nano\second}, avec un coefficient de variation robuste d'environ 14~\%.
\item Environ 18~\% (\qty{245}{\nano\second} contre 208)~: l'attente introduite de \qty{11.4}{\nano\second} a coûté plus que sa propre durée,
car les lectures d'horodatage de l'attente prennent aussi du temps.
\item Il est fixé par le 1~\% le plus lent des événements du passage, ceux touchés par les interruptions et les travaux propres au système, si bien qu'il dépend de
ce que le reste de la machine a fait pendant ce passage.
\item Son écart varie d'environ 20~\% d'une paire à l'autre, soit quatre fois la hausse à détecter.
\item Les médianes ignorent les passages gâchés par une rafale d'autres travaux, ce qui tirerait la moyenne.
\item Que, sans écart réel, les étiquettes A et B pourraient être inversées (l'entrelacement rend cela plausible)~; non que
les passages soient normaux, ni que A et B aient la même dispersion.
\item De 3 à 9~\% sur cent barrières, autour du seuil du test de 5~\% (cent barrières estiment un taux de 5~\% à quelques
points près)~; la tolérance supprime les hausses significatives inférieures à 2~\%.
\item 40.
\item 57~\%, 73~\% et 84~\%.
\item Environ 200 passages par côté, où le taux atteint 91~\%.
\item Environ 220, à peu près juste~: l'approximation suppose des passages normaux et utilise pour eux l'efficacité de la médiane~; les passages mesurés
sont concentrés avec quelques valeurs aberrantes éloignées, pour lesquelles la médiane fait un peu mieux que cela.
\item Environ $2 \times 200 \times \qty{30}{\milli\second} \approx \qty{12}{\second}$.
\item \textbf{Résultat retenu.} Sur cet ordinateur portable, dont les 99\up{e}~percentiles des passages varient d'environ 14~\% même au repos,
la barrière exige environ 200 passages entrelacés par côté (environ 12 secondes de benchmark) pour détecter une hausse de 5~\% avec une puissance de 90~\%, à un
taux de fausses alertes proche du seuil de 5~\% du test (l'approximation normale donne environ 220)~; une régression de 18~\% exige 40.
\item Moins de variation d'un passage à l'autre~; le nombre de passages nécessaires décroît comme le carré de la dispersion, jusqu'à environ 27 pour une dispersion de 5~\%.
\item Avec assez de passages, tout changement devient significatif~; la tolérance indique quels changements comptent.
\item Bien moins de passages, car les médianes sont stables~; mais elle manque ce que capte le 99\up{e}~percentile~: un chemin lent rare, comme une
allocation occasionnelle ou une ligne de journalisation sur une branche.
\item Le lancer sur les candidats à la fusion et chaque nuit, sur une machine dédiée, avec un nombre de passages fixé d'après le
bruit mesuré de la machine~; conserver les tests fonctionnels à chaque modification.
\item Les changements du réseau, du noyau, des interactions avec d'autres processus, du matériel, et les données qui diffèrent de
celles du jour du benchmark.
\item Recalculer le verdict et l'analyse de puissance avec d'autres statistiques ou seuils, et comparer des machines.
\item Assez de passages entrelacés pour le bruit mesuré de la machine, jugé par un test peu exigeant, et une tolérance.
\end{enumerate}
\end{solution}

\begin{solution}{iq:ll:testing-and-deploying-low-latency-systems:1}
Un test qui génère des entrées aléatoires et vérifie une propriété pour toutes. Pour un carnet d'ordres~: après toute suite d'ajouts,
de modifications et d'annulations, le meilleur cours acheteur est inférieur au meilleur cours vendeur, la quantité à chaque niveau est égale à la somme de ses ordres, et
la suppression de chaque ordre laisse un carnet vide.
\iqlookfor{une propriété, un générateur et la réduction des contre-exemples.}
\end{solution}

\begin{solution}{iq:ll:testing-and-deploying-low-latency-systems:2}
Muter les paquets enregistrés (octets, longueurs, numéros de séquence, troncatures), les injecter dans des tampons de taille exacte, avec les sanitizers, et
vérifier des invariants au-delà des plantages~: chaque message accepté tient dans son tampon, les numéros de séquence sont cohérents, le carnet reste valide après
chaque message, et la concordance avec un analyseur de référence.
\iqlookfor{sanitizers, invariants et référence différentielle.}
\end{solution}

\begin{solution}{iq:ll:testing-and-deploying-low-latency-systems:3}
Demander comment la mesure a été faite~: passages entrelacés des deux constructions sur la même machine, en nombre suffisant pour son bruit, un test avec un
intervalle de confiance. Relancer avec davantage de passages si besoin~; 3~\% peut être dans le bruit, ou sous la tolérance même s'il est réel.
\iqlookfor{entrelacement, bruit, un test et une tolérance.}
\end{solution}

\begin{solution}{iq:ll:testing-and-deploying-low-latency-systems:4}
Que les messages du client et sa gestion des réponses de la plateforme suivent le protocole, pas à pas, sur le
système de test de la plateforme. Ne sont pas couverts la stratégie de la firme, ses contrôles de risque au-delà de ce que le script exerce, sa latence ni son comportement sous charge
et face aux défaillances en production.
\iqlookfor{conformité au protocole, et non justesse ni performance.}
\end{solution}

\begin{solution}{iq:ll:testing-and-deploying-low-latency-systems:5}
Lancer un preflight de l'hôte~; exécuter la nouvelle construction en mode shadow sur des entrées réelles et comparer ses sorties~; la lancer sur une petite part du trafic (un canari),
avec des limites serrées, puis l'élargir~; revenir en arrière en redéployant la construction précédente, une procédure préparée et testée à l'avance, avec le
journal et les positions reportés.
\iqlookfor{mode shadow, canari, retour arrière préparé.}
\end{solution}

\begin{solution}{iq:ll:testing-and-deploying-low-latency-systems:6}
À chaque modification~: tests unitaires, \omterm{def:ll:testing-and-deploying-low-latency-systems:tests}{tests par propriétés} et \omterm{def:ll:testing-and-deploying-low-latency-systems:tests}{tests différentiels}, et \omterm{def:ll:testing-and-deploying-low-latency-systems:tests}{test par fuzzing} avec les sanitizers. Sur les candidats~: scénarios du simulateur et
rejeux des journées enregistrées, scripts de \omterm{def:ll:testing-and-deploying-low-latency-systems:cert}{certification par la bourse}, et la \omterm{def:ll:testing-and-deploying-low-latency-systems:gate}{barrière de régression de performance} sur une machine dédiée. À la mise en production~: preflight de l'hôte,
exécution en mode shadow, canari, déploiement progressif, retour arrière préparé. Chaque étape détecte ce que la précédente ne peut pas voir.
\iqlookfor{la pyramide, et ce que détecte chaque couche.}
\end{solution}
